\chapter{CONCLUSION AND FUTURE SCOPE}
\graphicspath{{Chapter5/}}

\section{Conclusion}
In this report, we addressed the challenging task of multi-keyword fuzzy search over encrypted cloud data \cite{wang2014privacy}.% [cite: 460]  
Although the symmetric cryptographic primitives had been successfully used for the basic keyword search, the exact keyword spellings and the traditional symmetric key distribution model had become the main obstacles for the real-world application of the symmetric primitives \cite{song2000practical, li2010fuzzy}.% [cite: 23, 24, 26, 438]  
The baseline analysis was conducted on the new paradigm, which employed the LSH functions with the assistance of the Bloom filters for the development of the secure document indexes \cite{wang2014privacy, bloom1970space}.% [cite: 36, 111]  
By utilizing the Euclidean distance for the modeling of the keyword similarity based on the bigram vectors, the baseline scheme enabled the multi-keyword fuzzy matching without the exponential increase in the size of the index \cite{wang2014privacy, datar2004locality}.% [cite: 35, 37] 

However, to mitigate the inherent false negative limitations of conventional LSH, as well as the key management limitations of symmetric cryptography, this project has proposed a significantly enhanced architecture. By migrating the underlying matching algorithm from a set of symmetric matrix operations to a set of Inner Product Predicate Encryption (IPE) operations, the proposed scheme enables data owners to securely outsource data via a universally accessible public key, as well as dynamically delegate search privileges to users via a master secret \cite{shen2009predicate}.

In addition, the integration of MP-LSH significantly enhances the overall accuracy of data retrieval by systematically examining adjacent, highly probable hash buckets, thereby ensuring that relevant data, even in the presence of minor typographical errors, are not discarded \cite{lv2007multi}. Ultimately, the proposed scheme provides a highly accurate, scalable, and secure data search environment.

\section{Future Scope}
While the proposed Asymmetric MP-LSH scheme significantly advances the state-of-the-art in secure data retrieval, several avenues remain open for future research and optimization:

\begin{itemize}
    \item \textbf{Integration of TF-IDF for Relevance Ranking:} The current inner product evaluation approach computes the raw count of matching LSH bits \cite{wang2014privacy}.% [cite: 209] 
    The future iterations can include the TF-IDF weights into the secure representations. This would enable the Cloud Server to rank the results based on the relevance of the documents and the significance of the keywords.
    
    \item \textbf{Optimized Dynamic Data Updates:} Although the existing approach eliminates the need to maintain a predefined global dictionary \cite{wang2014privacy},% [cite: 213, 218]  
    in the future, the focus should be on developing a highly optimized protocol for performing dynamic data updates, including insertions, deletions, and modifications, in the context of the Asymmetric IPE scheme. The aim is to perform these operations with low computational overhead without decrypting or re-indexing the surrounding untouched data.% [cite: 219]
    
    \item \textbf{Semantic and Natural Language Search:} While the current bigram vectorization approach very elegantly solves the typographical error problem, it does nothing for the ability to actually understand the concepts. Expanding the system to incorporate the capability for real semantic search, via the incorporation of secure word embeddings (such as Word2Vec), would be a significant step forward. This would allow the system to search for synonymous concepts (such as "automobile" and "car") directly on the encrypted data.
    
    \item \textbf{Hardware-Accelerated Cryptographic Evaluation:} The computational complexity of the bilinear pairing operations for evaluating the IPE for large data sets is a concern. As a future direction, hardware acceleration for the Cloud Server-side, possibly with GPUs or other hardware accelerators, might be explored for significantly reducing the search latency for large data sets.
\end{itemize}